\begin{derivation}[从\eqnrefs{eq:soft-external-leg-kernel-dp68}到\eqnrefs{eq:weinberg-soft-photon-factor-dp10}]
由 $p_i^2=m_i^2c^2$ 与 $k^2=0$，分母为

\begin{equation*}
 (p_i+k)^2-m_i^2c^2=2p_i\!\cdot k.
\end{equation*}
分子中的质量壳奇异部分由 Clifford 代数换序分离：
\begin{align*}
 \bar u_i\slashed\varepsilon^*(\slashed p_i+\slashed k+m_ic)
 &=\bar u_i\slashed\varepsilon^*\slashed p_i
   +\bar u_i\slashed\varepsilon^*\slashed k
   +m_ic\,\bar u_i\slashed\varepsilon^*,\\
 \bar u_i\slashed\varepsilon^*\slashed p_i
 &=2p_i\!\cdot\varepsilon^*\,\bar u_i
   -\bar u_i\slashed p_i\slashed\varepsilon^*.
\end{align*}
再用在壳 Dirac 方程 $\bar u_i\slashed p_i=m_ic\,\bar u_i$，质量项相消，留下
\begin{equation*}
 \bar u_i\slashed\varepsilon^*(\slashed p_i+\slashed k+m_ic)
 =2p_i\!\cdot\varepsilon^*\,\bar u_i
 +\bar u_i\slashed\varepsilon^*\slashed k.
\end{equation*}
最后一项为 $\order(\omega)$，除以 $p_i\!\cdot k=\order(\omega)$ 后只贡献 $\order(\omega^0)$；领先项因此为
\begin{equation*}
 g_{\rm em}Q_i
 \frac{p_i\!\cdot\varepsilon_\lambda^*}{p_i\!\cdot k}\,\bar u_i.
\end{equation*}
入射外腿的相邻传播子动量为 $p_i-k$，故
\begin{equation*}
 (p_i-k)^2-m_i^2c^2=-2p_i\!\cdot k,
\end{equation*}
从而相对多一个负号。对反费米子外腿，顶角的物理电荷本身已经带有与相应费米子相反的符号；其外线旋量满足
\[
 (\slashed p_i+m_ic)v(p_i)=0,
 \qquad
 \bar v(p_i)(\slashed p_i+m_ic)=0.
\]
将这两个在壳关系代入与上面相同的分子换序，质量项仍然抵消，领先分子仍化为 $2p_i\!\cdot\varepsilon^*$。因此外腿方向只通过传播子分母的符号进入，粒子种类的电荷符号则由 $Q_i$ 本身进入。定义 $\sigma_i^{\rm io}=+1$ 表示出射外腿、$\sigma_i^{\rm io}=-1$ 表示入射外腿后，每条带电外腿统一贡献
\[
 g_{\rm em}\,\sigma_i^{\rm io}Q_i
 \frac{p_i\!\cdot\varepsilon_\lambda^*}{p_i\!\cdot k}.
\]
对全部外腿求和即得\eqnrefs{eq:weinberg-soft-photon-factor-dp10}。

规范位移 $\varepsilon_\lambda^\mu\mapsto\varepsilon_\lambda^\mu+\zeta k^\mu$ 使软因子改变
\begin{equation*}
 \delta S=g_{\rm em}\zeta^*\sum_i\sigma_i^{\rm io}Q_i.
\end{equation*}
硬过程的电荷守恒给 $\sum_i\sigma_i^{\rm io}Q_i=0$，所以 $\delta S=0$。这同时检验了领先软因子的规范不变性。
\end{derivation}

\begin{derivation}[从\eqnrefs{eq:weinberg-soft-photon-factor-dp10}到\eqnrefs{eq:soft-radiator-definition-dp10}]
先把\eqnrefs{eq:weinberg-soft-photon-factor-dp10} 写成速度形式。对第 $i$ 条外腿，

\begin{equation*}
 p_i^\mu=\frac{E_i}{c}(1,\boldsymbol\beta_i),
 \qquad
 k^\mu=\frac{\hbar\omega}{c}(1,\mathbf n),
 \qquad
 |\mathbf n|=1,
\end{equation*}
所以
\begin{equation*}
 p_i\!\cdot k
 =\frac{E_\ii\hbar\omega}{c^2}
 (1-\mathbf n\!\cdot\boldsymbol\beta_i).
\end{equation*}
对物理横向偏振求和时，只有速度相对于 $\mathbf n$ 的横向分量
\begin{equation*}
 \boldsymbol\beta_i^\perp
 =\boldsymbol\beta_i-\mathbf n(\mathbf n\!\cdot\boldsymbol\beta_i)
\end{equation*}
保留下来。使用采用上述协变态外线归一化的一光子相空间因子，并代入 $g_{\rm em}^2=4\pi\alpha$，真实软辐射相对于硬过程的领先概率为
\begin{equation*}
 \dd\mathcal R
 =\frac{\alpha}{4\pi^2}
 \frac{\dd\omega}{\omega}\,\dd\Omega
 \left|
 \sum_i\sigma_i^{\rm io}Q_i
 \frac{\boldsymbol\beta_i^\perp}
 {1-\mathbf n\!\cdot\boldsymbol\beta_i}
 \right|^2.
\end{equation*}
该式的量纲检查是直接的：$\alpha$、$\dd\omega/\omega$、$\dd\Omega$ 与速度比均无量纲。对
$\lambda_{\rm IR}/\hbar<\omega<\Delta E_{\rm det}/\hbar$ 积分即得\eqnrefs{eq:soft-radiator-definition-dp10}，其中径向积分为
\begin{equation*}
 \int_{\lambda_{\rm IR}/\hbar}^{\Delta E_{\rm det}/\hbar}
 \frac{\dd\omega}{\omega}
 =\ln\frac{\Delta E_{\rm det}}{\lambda_{\rm IR}}.
\end{equation*}
\end{derivation}

\begin{derivation}[从\eqnrefs{eq:soft-radiator-definition-dp10}到\eqnrefs{eq:virtual-real-ir-relation-dp10}]
虚软光子连接外腿 $i,j$ 时，相邻费米子传播子的软极限为

\begin{equation*}
 \frac{\slashed p_i+\sigma_i^{\rm io}\slashed k+m_ic}
 {(p_i+\sigma_i^{\rm io} k)^2-m_i^2c^2+\ii0^+}
 =
 \frac{\slashed p_i+m_ic+\order(k)}
 {2\sigma_i^{\rm io} p_i\!\cdot k+\ii0^+}.
\end{equation*}
利用外腿 Dirac 方程把顶角两侧的分子约化为 $2p_i^\mu$，虚软修正的红外部分写成
\begin{equation*}
 \delta_{\rm virt}^{\rm soft}\mathcal M_0
 =-\frac12\mathcal M_0
 \sum_{ij}g_{\rm em}^2\sigma_i^{\rm io}\sigma_j^{\rm io}Q_iQ_j
 \int_{\rm soft}\frac{\dd^4 k}{(2\pi)^4}
 \frac{p_i\!\cdot p_j}
 {(k^2+\ii0^+)(p_i\!\cdot k+\ii0_i)(-p_j\!\cdot k+\ii0_j)}
 +\text{红外有限项}.
\end{equation*}
这里 $\ii0_i$ 记录外腿属于入射还是出射，决定相应的轮廓边界值。

对虚光子传播子取不连续度，
\begin{equation*}
 \frac1{k^2+\ii0^+}
 -\frac1{k^2-\ii0^+}
 =-2\pi\ii\,\delta(k^2).
\end{equation*}
正能割线进一步取 $\delta_+(k^2)=\Theta(k^0)\delta(k^2)$，于是四维圈积分中的虚光子线变成真实一光子相空间，并保留与 $\dd\mathcal R$ 相同的 $p_i,p_j,Q_iQ_j$ 成对核。光学定理把这部分解释为从无真实软光子通道流向真实软光子通道的概率，因此无辐射概率获得相反号。由此得到\eqnrefs{eq:virtual-real-ir-relation-dp10}。
\end{derivation}

\begin{derivation}[从\eqnrefs{eq:virtual-real-ir-relation-dp10}到\eqnrefs{eq:bloch-nordsieck-cancellation-dp10}]
先把正文软辐射积分中的角变量完全积分掉，定义只依赖硬外腿运动学的非负系数
\begin{equation*}
 \mathcal B
 \equiv\frac{\alpha}{4\pi^2}
 \int\dd\Omega\,
 \left|
 \sum_i\sigma_i^{\rm io}Q_i
 \frac{\boldsymbol\beta_i^\perp}
 {1-\mathbf n\cdot\boldsymbol\beta_i}
 \right|^2\ge0.
\end{equation*}
于是未分辨真实光子在
$\lambda_{\rm IR}<\hbar\omega<\Delta E_{\rm det}$ 内的概率权重为
\begin{equation*}
 \mathcal R(\lambda_{\rm IR},\Delta E_{\rm det})
 =\mathcal B\int_{\lambda_{\rm IR}/\hbar}^{\Delta E_{\rm det}/\hbar}
 \frac{\dd\omega}{\omega}
 =\mathcal B\ln\frac{\Delta E_{\rm det}}{\lambda_{\rm IR}}.
\end{equation*}
这里 $\lambda_{\rm IR}$ 是计算用调节量，而 $\Delta E_{\rm det}$ 是实验把“有光子”和“无可分辨光子”划开的真实分辨率，两者物理地位不同。

为了把虚圈中的同一对数看清，引入任意软--硬分隔能量 $E_s$，满足所用软展开在
$\hbar\omega<E_s$ 内成立。由前一推导中的割线关系，虚软圈的红外部分具有同一 $\mathcal B\,\dd\omega/\omega$ 核而符号相反，因此可写成
\begin{equation*}
 2\Re\frac{\delta_{\rm virt}^{\rm soft}\mathcal M_0}{\mathcal M_0}
 =-\mathcal B\ln\frac{E_s}{\lambda_{\rm IR}}
 +C_{\rm virt}(E_s),
\end{equation*}
其中 $C_{\rm virt}(E_s)$ 在 $\lambda_{\rm IR}\to0$ 时有限。把对数拆成
\begin{align*}
 -\mathcal B\ln\frac{E_s}{\lambda_{\rm IR}}
 &=-\mathcal B\ln\frac{\Delta E_{\rm det}}{\lambda_{\rm IR}}
   -\mathcal B\ln\frac{E_s}{\Delta E_{\rm det}},
\end{align*}
就得到正文\eqnrefs{eq:virtual-real-ir-relation-dp10}，其中
\[
 C_{\rm IR\,finite}
 =C_{\rm virt}(E_s)-\mathcal B\ln\frac{E_s}{\Delta E_{\rm det}}.
\]

现在展开到 $O(\alpha)$。无真实可分辨光子的虚修正概率为
\begin{align*}
 |\mathcal M_0+\delta_{\rm virt}\mathcal M_0|^2
 &=|\mathcal M_0|^2
 \left[1+2\Re\frac{\delta_{\rm virt}\mathcal M_0}{\mathcal M_0}\right]
 +O(\alpha^2),
\end{align*}
而一个未分辨真实软光子的概率为
$|\mathcal M_0|^2\mathcal R+O(\alpha^2)$。两类末态在 Fock 空间正交，所以包容概率相加而不是振幅相加：
\begin{align*}
 P_{\rm incl}
 &=|\mathcal M_0|^2
 \left[
 1-\mathcal R+C_{\rm IR\,finite}+\mathcal R
 \right]+O(\alpha^2)\\
 &=|\mathcal M_0|^2
 \left[1+C_{\rm IR\,finite}\right]+O(\alpha^2).
\end{align*}
因此所有 $\ln\lambda_{\rm IR}$ 严格抵消，取
$\lambda_{\rm IR}\to0$ 得到正文\eqnrefs{eq:bloch-nordsieck-cancellation-dp10}。留下的
$\Delta E_{\rm det}$ 依赖并不应消失：改变探测能量分辨率就改变了“哪些软光子被计入同一个包容事件”的测量定义。
\end{derivation}

\begin{derivation}[从\eqnrefs{eq:soft-independent-n-emission-dp68}到\eqnrefs{eq:bloch-nordsieck-exponentiation-dp10}]
领先软极限中，不同软光子从同一个经典化外腿电流独立发射。若单个软窗口的积分为 $\mathcal R$，则 $n$ 个不可分辨光子的组合因子为 $1/n!$，所以

\begin{equation*}
 P_n^{\rm soft}=\frac{\mathcal R^n}{n!}P_0^{\rm soft}.
\end{equation*}
幺正性要求所有不可分辨光子数的概率和为一：
\begin{align*}
 1
 &=\sum_{n=0}^{\infty}P_n^{\rm soft}
 =P_0^{\rm soft}\sum_{n=0}^{\infty}\frac{\mathcal R^n}{n!},\\
 &=P_0^{\rm soft}\ee^{\mathcal R}.
\end{align*}
因此
\begin{equation*}
 P_0^{\rm soft}=\ee^{-\mathcal R},
 \qquad
 P_n^{\rm soft}=\ee^{-\mathcal R}\frac{\mathcal R^n}{n!},
\end{equation*}
即\eqnrefs{eq:bloch-nordsieck-exponentiation-dp10}。若进一步限制多光子的总软能量而不是逐光子能量，独立 Poisson 形式会获得次领先相关；这属于领先软近似之外的修正。
\end{derivation}

\begin{derivation}[从\eqnrefs{eq:soft-collinear-dot-exact-dp68}到\eqnrefs{eq:soft-collinear-denominator-dp10}]
取电子四动量和光子四动量

\begin{equation*}
 p^\mu=(E/c,\mathbf p),
 \qquad
 k^\mu=(\hbar\omega/c,\hbar\omega\mathbf n/c),
\end{equation*}
其中 $|\mathbf n|=1$，$\mathbf p\cdot\mathbf n=p\cos\theta$。于是
\begin{equation*}
 p\!\cdot k
 =\frac{E\hbar\omega}{c^2}
 -\frac{p\hbar\omega}{c}\cos\theta
 =\frac{E\hbar\omega}{c^2}(1-\beta_{\rm rel}\cos\theta),
\end{equation*}
其中 $\beta_{\rm rel}=pc/E$。小角时
\begin{equation*}
 \cos\theta=1-\frac{\theta^2}{2}+\order(\theta^4).
\end{equation*}
质量壳关系给
\begin{equation*}
 \beta_{\rm rel}^2=1-\frac{m_{\mathrm e}^2c^4}{E^2},
\end{equation*}
令 $x\equiv m_{\mathrm e}^2c^4/E^2$。质量壳关系给 $\beta_{\rm rel}=\sqrt{1-x}$，所以
\begin{align*}
1-\beta_{\rm rel}
&=\frac{x}{1+\sqrt{1-x}}\\
&=\frac{x}{2}+O(x^2)
=\frac{m_{\mathrm e}^2c^4}{2E^2}
+O\!\left(\frac{m_{\mathrm e}^4c^8}{E^4}\right).
\end{align*}
保留 $\theta^2$ 和 $m_{\mathrm e}^2c^4/E^2$ 的最低非零阶，得到\eqnrefs{eq:soft-collinear-denominator-dp10}。这也显示电子质量如何把严格共线分母从零抬升为有限值。
\end{derivation}
